\chapter{Metodología y rigor}
\label{chap:metodologia}

\section{Metodología de trabajo}

Se aplicará un enfoque ágil, iterativo e incremental, adaptado al desarrollo
individual de este TFG. El primer ciclo establecerá las bases del proyecto y
producirá un producto mínimo viable (MVP) del flujo experimental: preparación
de los datos, ejecución de una línea base y obtención de predicciones y
métricas. Los ciclos posteriores incorporarán o ajustarán modelos a partir de
los resultados y las conclusiones de la iteración anterior. Cada incremento se
contrastará con los objetivos definidos y el proceso se repetirá hasta cubrir
satisfactoriamente el alcance. Este enfoque encaja con el carácter experimental
del trabajo, pues la evaluación de los modelos orientará las decisiones
posteriores.

La planificación se gestionará mediante un tablero Kanban en Notion. Cada tarea
indicará su objetivo, prioridad y criterio de finalización, y avanzará por los
estados pendiente, en curso y terminada. Al cerrar cada iteración se revisará el
trabajo completado y se priorizarán las tareas siguientes. Git mantendrá la
trazabilidad de los cambios en el código y en la memoria.

El código se desarrollará de forma modular, separando las responsabilidades de
preparación de datos, definición de modelos, ejecución de experimentos y
evaluación. Se seguirán principios de arquitectura limpia para facilitar la
incorporación y comparación de componentes durante las iteraciones.

\section{Seguimiento y rigor}

Se probarán la preparación de datos y las métricas, se conservarán
configuraciones, predicciones y resultados por experimento. Estos registros y
las tareas permitirán verificar las entregas y el cumplimiento de cada
objetivo: predicción de las cinco variables, comparación con la línea base y
medición del coste en el entorno definido. El progreso se revisará con la
dirección según su disponibilidad. Los modelos se elegirán sin consultar la
prueba final; se documentarán fallos y límites de los datos.
